\documentclass[10pt]{article}
\usepackage[a4paper,top=1.9cm,bottom=1.9cm,left=2.3cm,right=2.3cm]{geometry}
\usepackage[T1]{fontenc}
\usepackage{newtxtext}
\usepackage{hyperref}
\hypersetup{colorlinks=true,urlcolor=blue!55!black}
\setlength{\parindent}{0pt}
\setlength{\parskip}{0.42em}
\pagestyle{empty}
\begin{document}

10 November 2026

\vspace{0.2em}
The Editors\\
\emph{Nature Computational Science}

\vspace{1.0em}
\textbf{Re: submission of ``Bridging protein and nucleic-acid pockets via shared physicochemical
representations for 3D molecular generation with NA2smol''}

\vspace{0.4em}
Dear Editors,

Structure-based molecular generation has become routine for protein pockets, but nucleic acids
have been left behind, for two reasons. Part of the gap is data: curated nucleic-acid--small-molecule
resources hold hundreds of non-redundant entries against the millions of protein complexes available
for protein training. Part of it is representation. Most pocket-conditioned generators encode a
pocket through amino-acid-specific features for which there is no nucleotide analogue, so a model
trained on proteins cannot simply be pointed at an RNA site.

We report NA2smol, a Bayesian flow network generator that addresses both constraints together.
The pocket representation is restricted to descriptors that proteins and nucleic acids share:
pharmacophore type, hybridisation state, van der Waals radius and partial charge. Structural
supervision comes from ligand-only electron-density pretraining, transferred to the
pocket-conditioned pathway by a teacher--student alignment loss, with a cross-view contrastive
term completing the objective. One model generates three-dimensional ligands for either class of
binding site.

We benchmark against six structure-based generators spanning the main generation paradigms:
MolCRAFT, Pocket2Mol and ED2Mol, together with DiffSBDD, GraphBP and FLAG. On the RNA test set,
NA2smol retained 239 molecules across 17 of 27 evaluable targets under the most stringent tier of
a four-stage screening cascade, 0.68\% of 35,358 valid generations, and the protein-trained
baselines retained far fewer. On retrospective ECFP4 recovery, NA2smol leads both the RNA and protein benchmarks among these generators. The
same output, carried through post-processing of the generated library, produced a
bis-benzothiazole that reduced normalised thioflavin T fluorescence in a concentration-dependent
manner, with an apparent half-maximal concentration of 22.4~$\mu$M in a KRAS promoter
G-quadruplex assay.

The paper reports where NA2smol leads and where it does not. On conformational strain it is
mid-pack rather than leading, and the strain table reports success separately from the pass rate
so that geometries which could not be evaluated are visible as such. The experimental endpoint
rests on one compound selected through an integrated screen, which is what a first prospective
test of a generative pipeline can offer.

We believe the work suits \emph{Nature Computational Science} because its contribution is a
computational design question rather than a single application: how to build a representation and
a training strategy that transfer across two chemically different target classes when labelled
structural data exist for only one. Nucleic-acid pockets are a demanding test of that question,
and the paper reports the pipeline end to end. The complete training code, generated-molecule
tables, screening summaries and figure source data are available at
\url{https://github.com/Gallo13th/NA2mol} (commit \texttt{bac48aa}), and the manuscript states
the target universes and yield denominators explicitly so that the reported counts can be
recomputed.

The manuscript is original, is not under consideration elsewhere, and all authors have approved
this submission. We declare no competing interests. We would be glad to suggest referees in
generative molecular design, nucleic-acid recognition and structure-based screening.

Thank you for considering our work.

\vspace{0.8em}
Yours sincerely,

\vspace{1.4em}
Zhi John Lu\\
Associate Professor, School of Life Sciences, Tsinghua University\\
on behalf of all authors\\
\href{mailto:zhilu@tsinghua.edu.cn}{zhilu@tsinghua.edu.cn}

\end{document}
